\documentclass[12pt]{article}
\usepackage{amsmath, amssymb}
\usepackage[margin=1in]{geometry}
\setlength{\parskip}{6pt}

\title{QUBO Formulation: Financial Fraud Feature Selection}
\date{}

\begin{document}
\maketitle

\subsection*{Financial Fraud Feature Selection}

\paragraph{Problem Statement.}
Given a set of $N$ candidate features for fraud detection, each feature $i \in \{0, \dots, N-1\}$ is associated with a predictive value $v_i \geq 0$. Each unordered pair of distinct features $(i,j)$ with $i < j$ carries a redundancy penalty $r_{ij} \geq 0$. The task is to select exactly $K$ features such that the total predictive value is maximized while subtracting the pairwise redundancy penalties among the selected features.

\paragraph{Decision Variables.}
For each candidate feature $i \in \{0, \dots, N-1\}$, introduce a binary decision variable $x_i \in \{0,1\}$, where $x_i = 1$ indicates that feature $i$ is included in the selected set, and $x_i = 0$ otherwise.

\paragraph{QUBO Derivation.}
\textbf{Step 1.} The original problem is a maximization. To cast it as a QUBO minimization, we negate the objective function:
\begin{align}
\min_{x \in \{0,1\}^N} \left( -\sum_{i=0}^{N-1} v_i x_i + \sum_{0 \leq i < j \leq N-1} r_{ij} x_i x_j \right).
\end{align}

\textbf{Step 2.} The cardinality constraint requires exactly $K$ features to be selected:
\begin{align}
\sum_{i=0}^{N-1} x_i = K.
\end{align}

\textbf{Step 3.} We enforce the constraint via a quadratic penalty term with weight $\lambda > 0$.
\begin{align}
\text{(a) Squared penalty:} \quad & \lambda \left( \sum_{i=0}^{N-1} x_i - K \right)^2. \\
\text{(b) Expansion:} \quad & \lambda \left( \sum_{i=0}^{N-1} x_i^2 - 2K \sum_{i=0}^{N-1} x_i + K^2 + \sum_{i \neq j} x_i x_j \right). \\
& \text{Using the idempotent property } x_i^2 = x_i \text{ and noting } \sum_{i \neq j} x_i x_j = 2 \sum_{0 \leq i < j \leq N-1} x_i x_j, \text{ we obtain:} \\
& \lambda \left( \sum_{i=0}^{N-1} (1 - 2K) x_i + 2 \sum_{0 \leq i < j \leq N-1} x_i x_j + K^2 \right). \\
\text{(c) Penalty contribution to } H: \quad & \lambda \sum_{i=0}^{N-1} (1 - 2K) x_i + 2\lambda \sum_{0 \leq i < j \leq N-1} x_i x_j + \lambda K^2.
\end{align}

\textbf{Step 4.} Combining the negated objective from Step 1 with the penalty contribution from Step 3 yields the full QUBO Hamiltonian $H(x)$:
\begin{align}
H(x) = \sum_{i=0}^{N-1} \left( -v_i + \lambda(1 - 2K) \right) x_i + \sum_{0 \leq i < j \leq N-1} \left( r_{ij} + 2\lambda \right) x_i x_j + \lambda K^2.
\end{align}

\textbf{Step 5.} Reading off the coefficients from the general form $H(x) = \sum_{i} Q_{i,i} x_i + \sum_{i < j} Q_{i,j} x_i x_j + c$, we obtain the closed-form expressions for the upper-triangular Q matrix entries and the constant offset:
\begin{align}
Q_{i,i} &= -v_i + \lambda(1 - 2K), \quad \forall i \in \{0, \dots, N-1\}, \\
Q_{i,j} &= r_{ij} + 2\lambda, \quad \forall 0 \leq i < j \leq N-1, \\
c &= \lambda K^2.
\end{align}

\paragraph{Penalty Weight Justification.}
Let $V_{\max} = \max_{i} v_i$ denote the maximum predictive value of any single feature. The maximum possible objective gain achievable by selecting any subset of features is bounded by $N V_{\max}$. To guarantee that constraint violation is never energetically favorable, the penalty weight $\lambda$ must exceed the maximum possible improvement obtainable by an infeasible assignment. A sufficient condition is $\lambda > N V_{\max}$. In practice, choosing $\lambda$ significantly larger than any individual $v_i$ ensures that the energy cost of deviating from the cardinality constraint strictly dominates any potential gain from feature selection, thereby enforcing feasibility.

\paragraph{Correctness Argument.}
Minimizing $H(x)$ over $\{0,1\}^N$ recovers the optimal solution because (i) for any feasible assignment satisfying $\sum_{i=0}^{N-1} x_i = K$, the penalty term vanishes and $H(x)$ reduces exactly to the negated original objective; thus, the global minimum over the feasible subspace corresponds to the maximum of the original objective. (ii) For any infeasible assignment, the penalty term contributes a strictly positive energy $\lambda (\sum_{i=0}^{N-1} x_i - K)^2 \geq \lambda$. By selecting $\lambda$ sufficiently large, the energy of any infeasible point exceeds the energy of the optimal feasible point. Consequently, the global minimum of $H(x)$ must lie within the feasible region and correspond to the optimal feature set.

\paragraph{Illustrative Example.}
Consider a concrete instance with $N=3$ candidate features and $K=2$ required selections. Let the predictive values be $v_0=5, v_1=3, v_2=4$ and the redundancy penalties be $r_{01}=1, r_{02}=2, r_{12}=1$. Using $\lambda=10000$, the general formulas instantiate as follows.

The concrete objective terms are $-5x_0 - 3x_1 - 4x_2 + 1x_0x_1 + 2x_0x_2 + 1x_1x_2$.
The concrete penalty terms expand to $10000 \left( (1-4)(x_0+x_1+x_2) + 2(x_0x_1+x_0x_2+x_1x_2) + 4 \right)$.
Combining these yields the explicit Hamiltonian:
\begin{align}
H(x) &= (-5 - 20000)x_0 + (-3 - 20000)x_1 + (-4 - 20000)x_2 \nonumber \\
&\quad + (1 + 20000)x_0x_1 + (2 + 20000)x_0x_2 + (1 + 20000)x_1x_2 + 40000 \nonumber \\
&= -20005x_0 - 20003x_1 - 20004x_2 + 20001x_0x_1 + 20002x_0x_2 + 20001x_1x_2 + 40000.
\end{align}
The corresponding Q matrix entries are $Q_{0,0}=-20005$, $Q_{1,1}=-20003$, $Q_{2,2}=-20004$, $Q_{0,1}=20001$, $Q_{0,2}=20002$, $Q_{1,2}=20001$, with offset $c=40000$.

Evaluating all $\binom{3}{2}=3$ feasible assignments:
\begin{itemize}
    \item $x=(1,1,0)$: $H = -20005 - 20003 + 20001 + 40000 = 39993$.
    \item $x=(1,0,1)$: $H = -20005 - 20004 + 20002 + 40000 = 39993$.
    \item $x=(0,1,1)$: $H = -20003 - 20004 + 20001 + 40000 = 39994$.
\end{itemize}
The minimum energy is $39993$, achieved by $x=(1,1,0)$ and $x=(1,0,1)$. These correspond to feature sets $\{0,1\}$ and $\{0,2\}$, both yielding a maximum total score of $7$.

\end{document}